\section{Efecto sobre la distribución de precios del mercado local}
\label{anexo:mercado}

El cuerpo estima el efecto de la entrada sobre el precio de cada incumbente. Ese
ejercicio responde cuánto baja el precio, pero no dónde del mercado local se
produce la baja, y esa segunda pregunta es la que determina quién se beneficia.
Si la entrada arrastra hacia abajo tanto a la estación más cara como a la más
barata, la ganancia se reparte entre todos los consumidores; si comprime el
mercado desde abajo, se concentra en quienes comparan precios antes de comprar.

\subsection{Diseño}

La unidad de observación cambia: es el mercado y no la estación. Un mercado se
define como el círculo de radio dado alrededor de una incumbente focal, y las
variables de resultado son el precio máximo, el promedio y el mínimo entre todas
las estaciones activas dentro de él. Se estima la misma especificación de
\eqref{eq:did}, con efecto fijo de mercado en lugar de estación.

Dos decisiones merecen mención. La primera es que el precio de la entrante
\emph{sí} entra en los agregados desde el mes en que abre. \textcite{FISCHER2025}
excluyen a las entrantes del análisis a nivel de estación, pero advierten que sus
precios importan para la elección del consumidor y las incorporan justamente en
este ejercicio; el criterio se replica aquí. Omitirlas sesgaría el resultado
precisamente donde interesa, ya que en Chile la entrante típica es de bandera
blanca y tiende a fijar precios bajos.

La segunda es que se exige al mercado tener al menos dos estaciones, medidas
\emph{antes} del tratamiento. Si la restricción se aplicara con el conteo
contemporáneo, seleccionaría sobre la propia entrada, que es lo que altera ese
conteo.

Se excluye de este ejercicio el efecto fijo de marca por año que sí lleva la
especificación principal: el resultado pertenece al mercado y no a la estación
focal, de modo que la marca de esta última no es el control apropiado.

\subsection{Resultados}

La \Cref{fig:mercado93} y la \Cref{fig:mercadodi} presentan los tres agregados
para la gasolina de 93 octanos y para el diésel, en las dos definiciones de
mercado.

\begin{figure}[H]
\centering
\includegraphics[width=\textwidth]{mercado_max_media_min_p93.pdf}
\caption{Efecto de la entrada sobre el precio máximo, medio y mínimo del mercado: gasolina de 93 octanos}
\label{fig:mercado93}
\end{figure}

\begin{figure}[H]
\centering
\includegraphics[width=\textwidth]{mercado_max_media_min_pdi.pdf}
\caption{Efecto de la entrada sobre el precio máximo, medio y mínimo del mercado: diésel}
\label{fig:mercadodi}
\end{figure}

El patrón es inequívoco y se repite en ambos combustibles: la entrada comprime
la distribución de precios del mercado desde abajo. Dentro de un kilómetro, el
precio mínimo cae 5,5 pesos por litro en la gasolina de 93 octanos y 5,4 en el
diésel, mientras que el máximo no se mueve de manera distinguible de cero. El
promedio, que es el objeto que estima la especificación principal, se ubica entre
ambos.

La lectura económica es directa. La estación que ya era la más barata del mercado
responde a la entrada bajando aún más, mientras que la más cara no reacciona.
Esto es consistente con un mercado en que conviven consumidores con distinto
grado de información o distinto costo de búsqueda: quien compara precios se
beneficia de una caída varias veces mayor que quien compra en la estación más
próxima sin comparar. Es también coherente con los resultados de heterogeneidad
del cuerpo, aunque se trate de objetos distintos —allí los cuantiles son de la
distribución de márgenes entre estaciones, aquí de los estadísticos de orden
dentro de cada mercado local— y con lo que encuentran \textcite{FISCHER2025} para
Alemania.

Conviene señalar una limitación de precisión. El error estándar del efecto sobre
el precio máximo es grande, de modo que la afirmación de que el máximo no
responde es un no rechazo y no una estimación precisa de cero. Establecer con
propiedad que la entrada desplaza la distribución en el sentido de dominancia
estocástica de primer orden exigiría los contrastes formales sobre la
distribución completa que se discuten en las extensiones de las conclusiones, y
que la muestra chilena no permite sostener.
